% Generated by scripts/assemble.py; corrections belong in source/corrections.json.
\examhead{191}{西安电子科技大学 2025 年数学分析试卷}{180}
% p181-b19; source PDF page 181
\par\Needspace{5\baselineskip}
\blockanchor{p181-b19}
\majorquestion{一}计算题. 每题 7 分, 共 42 分.\par
% p181-b20; source PDF page 181
\blockanchor{p181-b20}
\question{1}$\lim_{x\to 0}\frac{1 - (\cos x)^{\sin x}}{x^3}.$\par
% p181-b21; source PDF page 181
\blockanchor{p181-b21}
\question{2}$\int \frac{1 + \arcsin x}{(1 - x^2)^{\frac{3}{2}}} \, \mathrm{d}x.$\par
% p181-b22; source PDF page 181
\blockanchor{p181-b22}
\question{3}设  $a_{n}=\int_{-1}^{1}\frac{x^{2n}}{1+e^{x}}\mathrm{d}x,\ n=0,1,2,\cdots$ ，求级数  $\sum_{n=0}^{\infty}(-1)^{n}a_{n}$  的和.\par
% p181-b23; source PDF page 181
\blockanchor{p181-b23}
\question{4}计算 $I = \iint_{\Sigma} \sqrt{x^2 + y^2} \, \mathrm{d}S$, 其中 $\Sigma: x^2 + y^2 + z^2 = 1, z \geq 0$.\par
% p181-b24; source PDF page 181
\blockanchor{p181-b24}
\question{5}计算 $\iint_{S} x^{3} \, \mathrm{d}y \, \mathrm{d}z$, $S$ 是椭球面 $\frac{x^{2}}{a^{2}} + \frac{y^{2}}{b^{2}} + \frac{z^{2}}{c^{2}} = 1$ 的上半部分, 取上侧.\par
% p182-b00; source PDF page 182
\blockanchor{p182-b00}
\question{6}设 $f(x)$ 在 $a$ 的某邻域内具有连续的 $n + 1$ 阶导函数, 求 $\lim_{x \to a} \frac{\mathrm{d}^n}{\mathrm{d}x^n} \left[ \frac{f(x) - f(a)}{x - a} \right]$.\par
% p182-b01; source PDF page 182
\par\Needspace{5\baselineskip}
\blockanchor{p182-b01}
\majorquestion{二}判断下列命题是否正确, 若正确请证明, 错误则给出反例. 每题 7 分, 共 28 分.\par
% p182-b02; source PDF page 182
\blockanchor{p182-b02}
\question{1}函数列 $\{f_n(x)\} = \left\{\frac{1}{1 + n^2x}\right\}$ 在 $(0,1)$ 上一致收敛.\par
% p182-b03; source PDF page 182
\blockanchor{p182-b03}
\question{2}对任意正整数 $n$ 和 $p$, 均有 $|x_{n+p} - x_n| < \frac{p}{n^2}$, 则数列 $\{x_n\}$ 是柯西列.\par
% p182-b04; source PDF page 182
\par\Needspace{6\baselineskip}
\blockanchor{p182-b04}
\question{3}设二元函数 $f(x, y)$ 有二阶连续偏导数，$(x_0, y_0)$ 是 $f(x, y)$ 的严格极小值点，则\par
\nopagebreak[4]
% p182-b05; source PDF page 182
\blockanchor{p182-b05}
\displayformula{f _ {x x} (x _ {0}, y _ {0}) f _ {y y} (x _ {0}, y _ {0}) - f _ {x y} ^ {2} (x _ {0}, y _ {0}) > 0.}
% p182-b06; source PDF page 182
\blockanchor{p182-b06}
\question{4}设级数 $\sum_{n=1}^{\infty} u_n$ 收敛, 则级数 $\sum_{n=1}^{\infty} u_n^3$ 收敛.\par
% p182-b07; source PDF page 182
\par\Needspace{5\baselineskip}
\par\Needspace{6\baselineskip}
\blockanchor{p182-b07}
\majorquestion{三}(13 分) 有 $p \in \mathbb{R}$, 讨论反常积分 $\int_{0}^{+\infty} \frac{\sin x}{x^p} \, \mathrm{d}x$ 的敛散性, 若收敛, 判断是条件收敛还是绝对收敛.
\majorquestion{四}(13 分) 设函数\par
\nopagebreak[4]
% p182-b08; source PDF page 182
\blockanchor{p182-b08}
\displayformula{k (x, y) = \left\{ \begin{array}{l l} y (1 - x), & y \leq x; \\ x (1 - y), & y > x. \end{array} \right.}
% p182-b09; source PDF page 182
\blockanchor{p182-b09}
$f(y)$ 在 $[0,1]$ 上连续. 证明: $u(x) = \int_{0}^{1} k(x,y)f(y)\mathrm{d}y$ 在 $[0,1]$ 上满足微分方程 $\frac{\mathrm{d}^2 u}{\mathrm{d}x^2} + f(x) = 0$.\par
% p182-b10; source PDF page 182
\par\Needspace{5\baselineskip}
\par\Needspace{6\baselineskip}
\blockanchor{p182-b10}
\majorquestion{五}（14分） $\varphi (y)$ 有连续的导函数，且对右半平面 $x > 0$ 内任意闭合曲线 $C$ ，均有\par
\nopagebreak[4]
% p182-b11; source PDF page 182
\blockanchor{p182-b11}
\displayformula{\oint_ {C} \frac {\varphi (y) \mathrm{d} x + 3 x y ^ {2} \mathrm{d} y}{2 x ^ {2} + y ^ {6}} = 0.}
% p182-b12; source PDF page 182
\blockanchor{p182-b12}
(1) 求  $\varphi(y)$  的表达式.\par
% p182-b13; source PDF page 182
\blockanchor{p182-b13}
(2) 设曲线 $L: x^2 + 2y^2 = 1$, 取逆时针方向, 求 $\oint_{L} \frac{\varphi(y) \mathrm{d}x + 3xy^2 \mathrm{d}y}{2x^2 + y^6}$.\par
% p182-b14; source PDF page 182
\par\Needspace{5\baselineskip}
\par\Needspace{6\baselineskip}
\blockanchor{p182-b14}
\majorquestion{六}(13 分) 设  $f_{i}(x)$  ( $i = 1, 2, \cdots, n$ ) 为 [a, b] 上的连续函数, 证明:\par
\nopagebreak[4]
% p182-b15; source PDF page 182
\blockanchor{p182-b15}
\displayformula{\sum_ {k = 1} ^ {n} \left(\int_ {a} ^ {b} f _ {k} (x) \mathrm{d} x\right) ^ {2} \leq \left(\int_ {a} ^ {b} \sqrt {\sum_ {k = 1} ^ {n} f _ {k} ^ {2} (x)} \mathrm{d} x\right) ^ {2}}
% p182-b16; source PDF page 182
\par\Needspace{5\baselineskip}
\blockanchor{p182-b16}
\majorquestion{七}(13 分) 设向量值函数 $f: (-1, 1) \to \mathbb{R}^2$ 有连续导数, 且 $f(0) = 0$, $f'(0) \neq 0$, 证明: 存在 $\varepsilon \in (0, 1)$, 使得 $\| f(t) \|$ 关于 $t$ 在 $(0, \varepsilon)$ 上单调递增.\par
% p182-b17; source PDF page 182
\par\Needspace{5\baselineskip}
\blockanchor{p182-b17}
\majorquestion{八}（14分）设函数 $f$ 在区间 $[a,b]$ 上有定义且处处存在极限，证明：\par
% p182-b18; source PDF page 182
\blockanchor{p182-b18}
(1) 对任意的 $\varepsilon > 0$, 在 $[a, b]$ 上使 $\left|\lim_{t \to x} f(t) - f(x)\right| > \varepsilon$ 的点 $x$ 至多有有限个.\par
% p182-b19; source PDF page 182
\blockanchor{p182-b19}
(2) $f$ 在 $[a, b]$ 上至多只有可列个间断点.\par
